% Luc Destombe --- licence CC BY-NC-SA 4.0
% https://creativecommons.org/licenses/by-nc-sa/4.0/deed.fr
% Fichier autonome : compiler avec pdflatex, deux fois (signets, nombre de pages).
\def\lyceeversion{1.0}
\documentclass[a4paper,10pt]{article}
\makeatletter
\newif\iflycee@niveaux
\newif\iflycee@cm
\newif\iflycee@palatino

\RequirePackage[utf8]{inputenc}
\RequirePackage[T1]{fontenc}
\RequirePackage[french]{babel}
\RequirePackage{etoolbox}

\newcommand{\ShorthandsOff}{\@ifundefined{shorthandoff}{}{\shorthandoff{;:!?}}}
\newcommand{\ShorthandsOn}{\@ifundefined{shorthandon}{}{\shorthandon{;:!?}}}
\ShorthandsOff
\AtBeginDocument{\ShorthandsOn}
\AtBeginEnvironment{tikzpicture}{\ShorthandsOff}

\RequirePackage{geometry}
\RequirePackage{fancyhdr}
\RequirePackage{lastpage}
\RequirePackage{multicol}
\RequirePackage{array}
\RequirePackage{enumitem}
\RequirePackage{xcolor}
\RequirePackage{needspace}

\RequirePackage{amsmath,amssymb}
\iflycee@cm\else\RequirePackage{fourier}\fi
\DeclareMathAlphabet{\mathcal}{OMS}{cmsy}{m}{n}
\renewcommand{\ttdefault}{lmtt}
\begingroup\catcode`\_=\active \catcode`\^=\active
\global\let_\sb \global\let^\sp
\endgroup
\AtBeginDocument{\catcode`\_=12 \mathcode`\_="8000
  \catcode`\^=12 \mathcode`\^="8000 }
\newlength{\retraitcalculs}
\setlength{\retraitcalculs}{2em}

\RequirePackage{tikz}
\usetikzlibrary{arrows.meta,calc,babel,shapes.misc}
\RequirePackage[most]{tcolorbox}
\tcbuselibrary{skins,breakable}

\definecolor{coulDef}{HTML}{1F4E79}
\definecolor{coulProp}{HTML}{7030A0}
\definecolor{coulRem}{HTML}{C55A11}
\definecolor{coulEx}{HTML}{2E7D32}
\definecolor{coulExo}{HTML}{B03A2E}
\definecolor{coulPart}{HTML}{1F4E79}
\definecolor{coulMeth}{HTML}{1B6E4F}
\definecolor{coulCourbe}{HTML}{0B5ED7}
\definecolor{coulCourbeB}{HTML}{C00000}
\definecolor{coulCourbeC}{HTML}{2E7D32}
\definecolor{coulGrille}{HTML}{8C8C8C}
\definecolor{coulRep}{HTML}{1B6E4F}
\definecolor{coulBar}{HTML}{2E7D32}

\newcommand{\R}{\mathbb{R}}
\newcommand{\Rs}{\mathbb{R}^{*}}
\newcommand{\N}{\mathbb{N}}
\newcommand{\Ns}{\mathbb{N}^{*}}
\newcommand{\Z}{\mathbb{Z}}
\newcommand{\Q}{\mathbb{Q}}
\renewcommand{\le}{\leqslant}
\renewcommand{\ge}{\geqslant}
\renewcommand{\approx}{\simeq}
\newcommand{\vect}[1]{\overrightarrow{#1}}
\newcommand{\norme}[1]{\left\lVert #1 \right\rVert}
\newcommand{\intervalle}[2]{\left[#1\,;#2\right]}
\RequirePackage{cancel}
\renewcommand{\CancelColor}{\color{coulExo}}
\newcommand{\fois}[1]{\times{\color{coulExo}#1}}
\newcommand{\barre}[1]{\cancel{#1}}
\newlength{\lycee@eqh}\newlength{\lycee@eqd}
\newsavebox{\lycee@eqa}\newsavebox{\lycee@eqb}\newsavebox{\lycee@eqc}
\newcommand{\lycee@eqcol}[1]{\hbox{$\renewcommand{\arraystretch}{1.35}%
  \begin{array}[t]{@{}>{\displaystyle}r@{}>{\displaystyle{}}l@{}}#1\end{array}$}}
\newcommand{\lycee@eqmesure}[1]{%
  \ifdim\ht#1>\lycee@eqh\setlength{\lycee@eqh}{\ht#1}\fi
  \ifdim\dp#1>\lycee@eqd\setlength{\lycee@eqd}{\dp#1}\fi}
\newcommand{\lycee@eqtrait}{\hspace{1.2em}\tikz[baseline=0pt]\draw[dash pattern=on 2pt off 2pt,
  line width=0.5pt] (0,\lycee@eqh) -- (0,-\lycee@eqd);\hspace{1.2em}}
\newcommand{\eqproduit}[2]{\par\smallskip\noindent C'est une équation produit :
  $A\times B=0 \iff A=0$ ou $B=0$.\par\nopagebreak\smallskip
  \begingroup\setlength{\lycee@eqh}{0pt}\setlength{\lycee@eqd}{0pt}%
  \setbox\lycee@eqa=\lycee@eqcol{#1}\setbox\lycee@eqb=\lycee@eqcol{#2}%
  \lycee@eqmesure\lycee@eqa\lycee@eqmesure\lycee@eqb%
  \noindent\hspace*{\retraitcalculs}\usebox\lycee@eqa\lycee@eqtrait\usebox\lycee@eqb\endgroup\par\smallskip}
\newcommand{\eqproduittrois}[3]{\par\smallskip\noindent C'est une équation produit :
  $A\times B\times C=0 \iff A=0$ ou $B=0$ ou $C=0$.\par\nopagebreak\smallskip
  \begingroup\setlength{\lycee@eqh}{0pt}\setlength{\lycee@eqd}{0pt}%
  \setbox\lycee@eqa=\lycee@eqcol{#1}\setbox\lycee@eqb=\lycee@eqcol{#2}\setbox\lycee@eqc=\lycee@eqcol{#3}%
  \lycee@eqmesure\lycee@eqa\lycee@eqmesure\lycee@eqb\lycee@eqmesure\lycee@eqc%
  \noindent\hspace*{\retraitcalculs}\usebox\lycee@eqa\lycee@eqtrait\usebox\lycee@eqb\lycee@eqtrait\usebox\lycee@eqc\endgroup
  \par\smallskip}

\tikzset{grille/.style={coulGrille,line width=0.4pt}}
\newcommand{\quadrillage}[4]{\draw[grille,step=1] (#1,#2) grid (#3,#4);}
\tikzset{
  croix/.style={cross out,draw,line width=0.6pt,inner sep=0pt,minimum size=4pt},
  croixdroite/.style={croix,rotate=45,line width=0.8pt},
}
\newcommand{\pt}[3]{\path (#1) node[croix]{} node[#2,font=\small,inner sep=2.5pt]{$#3$};}

\newcommand{\lycee@repere}[4]{%
  \draw[grille,step=1] (#1,#2) grid (#3,#4);
  \draw[-{Stealth[length=2mm]},thick] (#1,0)--({#3+0.4},0) node[below left=-1pt]{$x$};
  \draw[-{Stealth[length=2mm]},thick] (0,#2)--(0,{#4+0.4}) node[below left=-1pt]{$y$};
  \node[below left,font=\scriptsize,inner sep=1.5pt] at (0,0) {$0$};}
\newcommand{\lycee@gradx}[1]{\foreach \k/\l in {#1}{%
  \draw (\k,0.07)--(\k,-0.07) node[below,font=\scriptsize,inner sep=1.5pt]{$\l$};}}
\newcommand{\lycee@grady}[1]{\foreach \k/\l in {#1}{%
  \draw (0.07,\k)--(-0.07,\k) node[left,font=\scriptsize,inner sep=1.5pt]{$\l$};}}
\newcommand{\lycee@intff}[2]{\left[#1\,;#2\right]}
\newcommand{\lycee@intfo}[2]{\left[#1\,;#2\right[}
\newcommand{\lycee@intof}[2]{\left]#1\,;#2\right]}
\newcommand{\lycee@intoo}[2]{\left]#1\,;#2\right[}
\newcommand{\lycee@ens}[1]{\left\{#1\right\}}
\AtBeginDocument{%
  \providecommand{\repere}{\lycee@repere}%
  \providecommand{\gradx}{\lycee@gradx}%
  \providecommand{\grady}{\lycee@grady}%
  \providecommand{\Cf}{\mathcal{C}\sb{f}}%
  \providecommand{\Cg}{\mathcal{C}\sb{g}}%
  \providecommand{\intff}{\lycee@intff}%
  \providecommand{\intfo}{\lycee@intfo}%
  \providecommand{\intof}{\lycee@intof}%
  \providecommand{\intoo}{\lycee@intoo}%
  \providecommand{\ens}{\lycee@ens}}

\newlist{questions}{enumerate}{3}
\setlist[questions,1]{label=\arabic*\textdegree),leftmargin=*,itemsep=2pt,topsep=3pt}
\setlist[questions,2]{label=\alph*),leftmargin=*,itemsep=1pt,topsep=2pt}
\setlist[questions,3]{label=\roman*),leftmargin=*,itemsep=1pt,topsep=1pt}
\setlist[itemize]{label=\textbullet,leftmargin=*,itemsep=2pt,topsep=3pt}

\newcommand{\besoinplace}[1]{\par\penalty\@M\begingroup
  \dimen@=#1\relax\dimen@ii=\pagegoal\advance\dimen@ii-\pagetotal
  \ifdim\dimen@>\dimen@ii\ifdim\dimen@ii>\z@\vfil\fi\break\fi\endgroup}

\newcommand{\lycee@pied}{\thepage/\pageref{LastPage}}
\newcommand{\entete}[2]{%
  \pagestyle{fancy}\fancyhf{}%
  \fancyhead[L]{\footnotesize\color{coulPart}\textbf{#1}}%
  \fancyhead[R]{\footnotesize\color{coulPart}\textbf{#2}}%
  \fancyfoot[C]{\footnotesize\lycee@pied}%
  \renewcommand{\headrulewidth}{0.4pt}%
  \renewcommand{\headrule}{\hbox to\headwidth{%
    \color{coulPart!40}\leaders\hrule height \headrulewidth\hfill}}}

\RequirePackage{ccicons}
\newcommand{\lycee@licence}{{\scriptsize\color{gray}%
  \copyright\ Luc Destombe --- \ccbyncsa\ CC BY-NC-SA 4.0}}
\newif\iflycee@licence \lycee@licencetrue
\newcommand{\sanslicence}{\lycee@licencefalse}
\AtBeginDocument{\iflycee@licence\AddToHookNext{shipout/foreground}{%
  \put(0,0){\rlap{\kern\dimexpr1in+\hoffset+\oddsidemargin+\textwidth\relax
    \llap{\raisebox{-\dimexpr1in+\voffset+\topmargin+\headheight+\headsep
      +\textheight+\footskip\relax}{\lycee@licence}}}}}\fi}

\newcommand{\lycee@site}{faire-des-maths.fr}
\newcommand{\lycee@version}{{\scriptsize\color{gray}\lycee@site
  \ifdefined\lyceeversion\ --- Version~\lyceeversion\fi}}
\AtBeginDocument{\iflycee@licence\AddToHookNext{shipout/foreground}{%
  \put(0,0){\rlap{\kern\dimexpr1in+\hoffset+\oddsidemargin\relax
    \raisebox{-\dimexpr1in+\voffset+\topmargin+\headheight+\headsep
      +\textheight+\footskip\relax}{\lycee@version}}}}\fi}

\newcommand{\formulesserrees}{%
  \setlength{\abovedisplayskip}{3pt}\setlength{\belowdisplayskip}{3pt}%
  \setlength{\abovedisplayshortskip}{2pt}\setlength{\belowdisplayshortskip}{3pt}}

\setlength{\parindent}{0pt}

\RequirePackage[implicit=false,hidelinks,pdfencoding=auto,psdextra,bookmarksopen,
  bookmarksopenlevel=1,pdfcreator={},pdfproducer={}]{hyperref}
\RequirePackage{bookmark}
\hypersetup{pdfauthor={Luc Destombe},pdfkeywords={CC BY-NC-SA 4.0}}
\newcounter{lycee@signet}
\newcommand{\signet}[3][]{\stepcounter{lycee@signet}%
  \edef\lycee@dest{\if\relax\detokenize{#1}\relax
    signet.\arabic{lycee@signet}\else#1\fi}%
  \hypertarget{\lycee@dest}{}\bookmark[level=#2,dest=\lycee@dest]{#3}}
\pdfstringdefDisableCommands{%
  \def\R{R}\def\N{N}\def\Z{Z}\def\Q{Q}\def\vect#1{#1}\def\Cf{Cf}\def\Cg{Cg}\def\fois#1{×#1}%
  \def\barre#1{#1}\def\intervalle#1#2{[#1;#2]}\def\intff#1#2{[#1;#2]}%
  \def\textsuperscript#1{#1}\def\dfrac#1#2{#1/#2}\def\frac#1#2{#1/#2}%
  \def\vec#1{#1⃗}}%
\RequirePackage[official]{eurosym}

\geometry{top=0.8cm,bottom=1.3cm,left=1cm,right=1cm,footskip=0.6cm}

\pagestyle{fancy}
\fancyhf{}
\renewcommand{\headrulewidth}{0pt}
\fancyfoot[C]{\footnotesize Page \thepage{} / \pageref{LastPage}}

\newcommand{\titrefiche}[2]{%
  \begin{center}
    \Large\textbf{#1}\\[2pt]
    \large\textbf{#2}\\[2pt]
  \end{center}
  \vspace{0.10cm}}

\setlength{\columnseprule}{0.4pt}
\setlength{\columnsep}{15pt}
\renewcommand{\columnseprulecolor}{\color{black!45}}
\setlength{\multicolsep}{2pt}

\newcounter{partiefiche}
\renewcommand{\thepartiefiche}{\Roman{partiefiche}}
\newif\ifapresbandeau
\newcommand{\lycee@bandeau}[1]{%
  \par\penalty-600
  \vspace{0.08cm}%
  \noindent\colorbox{coulPart}{%
    \parbox{\dimexpr\linewidth-2\fboxsep\relax}{%
      \color{white}\bfseries\raggedright\strut#1\strut}}%
  \nopagebreak\par\nopagebreak
  \vspace{0.06cm}\nopagebreak
  \apresbandeautrue}
\newcommand{\partiefiche}[1]{%
  \refstepcounter{partiefiche}%
  \lycee@bandeau{\signet{1}{\thepartiefiche{} --- #1}\thepartiefiche{} --- #1}}
\newcommand{\partiefichesansnum}[1]{\lycee@bandeau{\signet{1}{#1}#1}}

\newcounter{exo}
\newlength{\espaceexo}\setlength{\espaceexo}{7pt}
\newlength{\espacetitre}\setlength{\espacetitre}{2pt}
\newcommand{\exo}[1][\relax]{%
  \par
  \ifapresbandeau\apresbandeaufalse\else\penalty-150\addvspace{\espaceexo}\fi
  \refstepcounter{exo}%
  \noindent\signet[exercice-\arabic{exo}]{2}{Exercice \arabic{exo}}%
  \textbf{\color{coulExo}\underline{Exercice \theexo}}%
  \ifx\relax#1\relax\else\textbf{ : #1}\fi
  \par\nopagebreak\vspace{\espacetitre}\nopagebreak
  \ignorespaces}

\setlist[enumerate]{nosep,leftmargin=*,itemsep=2pt,topsep=2pt}
\setlist[itemize]{nosep,leftmargin=*,topsep=2pt}
\newcommand{\choix}[3]{\par\hspace*{1em}%
  \makebox[0.3\linewidth][l]{a) #1}\makebox[0.3\linewidth][l]{b) #2}c) #3}
\newcommand{\deux}[2]{\par\hspace*{0.6em}%
  \makebox[0.48\linewidth][l]{#1}#2}
\newcommand{\trois}[3]{\par\hspace*{0.6em}%
  \makebox[0.32\linewidth][l]{#1}\makebox[0.32\linewidth][l]{#2}#3}

\renewenvironment{center}{\par\addvspace{3pt}\centering}{\par\addvspace{3pt}}
\formulesserrees
\AtBeginDocument{\formulesserrees}
\clubpenalty=10000
\widowpenalty=10000
\displaywidowpenalty=10000
\brokenpenalty=10000
\setlength{\parskip}{0pt}

\RequirePackage{ragged2e}
\setlength{\RaggedRightRightskip}{0pt plus 1fil}
\AtBeginDocument{\RaggedRight\binoppenalty=10000 \relpenalty=10000 }
\g@addto@macro\@parboxrestore{\RaggedRight}
\makeatother
\tcbuselibrary{listings}
\geometry{bottom=1cm}

\setlength{\espaceexo}{4pt}
\setlist[enumerate]{nosep,leftmargin=*,itemsep=1pt,topsep=1pt}

\newcommand{\Card}{\operatorname{Card}}

\newtcblisting{python}{%
  listing only,enhanced,
  colback=black!3,colframe=black!35,boxrule=0.5pt,arc=1pt,
  left=4pt,right=4pt,top=1pt,bottom=1pt,
  listing options={language=Python,basicstyle=\ttfamily\small,
    keywordstyle=\color{coulMeth}\bfseries,columns=fullflexible,
    keepspaces=true,showstringspaces=false,aboveskip=0pt,belowskip=0pt}}

\begin{document}

\titrefiche{Seconde --- Probabilités}{Fiche d'exercices du chapitre}

\begin{multicols}{2}\raggedcolumns

\partiefiche{Vocabulaire des probabilités}

\exo[Univers et événements]
On tire au hasard une boule dans une urne qui contient $20$ boules numérotées de $1$
à $20$, et on note son numéro.
\begin{enumerate}[label=\alph*)]
  \item Écrire l'univers $\Omega$, puis donner $\Card(\Omega)$.
  \item Écrire sous forme d'ensemble les événements $A$ : « le numéro est un
    multiple de $6$ » et $B$ : « le numéro est supérieur ou égal à $17$ ».
  \item Écrire l'événement $C$ : « le numéro est un multiple de $25$ ». Comment
    s'appelle un tel événement ?
  \item La boule n°$18$ sort. Lesquels des événements $A$, $B$ et $C$ sont
    réalisés ? Répondre par une phrase.
\end{enumerate}

\exo[Un univers de mots]
On écrit chacun des mots ci-dessous sur un carton, on place les cartons dans un sac
et on tire un carton au hasard :
\begin{center}
$\Omega=\ens{\text{lac}\,;\text{mer}\,;\text{rio}\,;\text{sol}\,;\text{ver}\,;
\text{mur}\,;\text{roc}\,;\text{lit}}$
\end{center}
\begin{enumerate}[label=\alph*)]
  \item Donner $\Card(\Omega)$.
  \item Écrire les événements $A$ : « le mot contient la lettre r » et $B$ : « le
    mot contient la lettre o ».
  \item Écrire l'événement $C$ : « le mot commence par une voyelle », puis compléter
    par une phrase : « L'événement $C$ est … ».
  \item Compléter par une phrase : « L'événement $D$ : “le mot a trois lettres” est
    … ».
\end{enumerate}

\exo[Deux lancers d'une pièce]
On lance deux fois de suite une pièce de monnaie et on note les faces obtenues dans
l'ordre : PF signifie « pile, puis face ».
\begin{enumerate}[label=\alph*)]
  \item Écrire l'univers $\Omega$.
  \item Écrire les événements $A$ : « on obtient au moins une fois pile » et $B$ :
    « on obtient deux fois la même face ».
  \item On obtient FF. Lesquels des événements $A$ et $B$ sont réalisés ?
\end{enumerate}

\partiefiche{Calcul d'une probabilité}

\exo[Un dé à vingt faces]
On lance un dé équilibré à $20$ faces numérotées de $1$ à $20$. Calculer les
probabilités des événements suivants.
\begin{enumerate}[label=\alph*)]
  \item $A$ : « obtenir un nombre pair ».
  \item $B$ : « obtenir un multiple de $6$ ».
  \item $C$ : « obtenir un multiple de $7$ ».
  \item $D$ : « obtenir un nombre inférieur ou égal à $20$ ».
\end{enumerate}

\exo[Des jetons numérotés]
Un sac contient $12$ jetons indiscernables au toucher : $5$ portent le numéro $1$,
$4$ le numéro $2$ et $3$ le numéro $3$. On tire un jeton au hasard et on note son
numéro $n$.
\begin{enumerate}[label=\alph*)]
  \item Expliquer pourquoi les issues de l'expérience sont les $12$ jetons, et non
    les $3$ numéros.
  \item Calculer les probabilités des événements $A$ : « $n$ est impair », $B$ :
    « $n\ge 2$ » et $C$ : « $n=4$ ».
\end{enumerate}

\exo[Un jeu de 52 cartes]
Un jeu de $52$ cartes compte $4$ couleurs (pique et trèfle, noires ; cœur et
carreau, rouges) de $13$ cartes chacune : as, $2$ à $10$, valet, dame et roi. On
tire une carte au hasard. Calculer la probabilité de tirer :
\begin{enumerate}[label=\alph*)]
  \item un trèfle ;
  \item une carte noire ;
  \item une figure (valet, dame ou roi) ;
  \item un as.
\end{enumerate}

\exo[Le sac le plus favorable]
Lina, Tom et Inès ont chacun un sac de billes et tirent une bille au hasard dans le
leur. Lina a $4$ billes rouges ; Tom a $6$ billes rouges et $18$ noires ; Inès a $50$
billes rouges et $2$ noires.
\begin{enumerate}[label=\alph*)]
  \item Calculer, pour chacun, la probabilité de tirer une bille rouge. Qui a le plus
    de chances ?
  \item Combien de billes noires faut-il ajouter dans le sac de Lina pour qu'elle ait
    la même probabilité que Tom de tirer une bille rouge ?
\end{enumerate}

\partiefiche{Intersection, réunion, événement contraire}

\exo[Une roue de loterie]
Une roue de loterie a $16$ secteurs de même taille, numérotés de $1$ à $16$. On la
fait tourner et on note le numéro obtenu. On note $A$ : « le numéro est un multiple
de $4$ » et $B$ : « le numéro est supérieur ou égal à $13$ ».
\begin{enumerate}[label=\alph*)]
  \item Écrire $A$, $B$, $A\cap B$, $A\cup B$ et $\overline{A}$ sous forme
    d'ensembles. $A$ et $B$ sont-ils incompatibles ?
  \item Calculer $P(A)$, $P(B)$ et $P(A\cap B)$.
  \item Calculer $P\big(\overline{A}\big)$ et $P(A\cup B)$ avec la propriété 3, puis
    vérifier en comptant les issues.
\end{enumerate}

\exo[Traduire par une phrase]
On choisit au hasard un spectateur d'un match de handball. On note $E$ : « le
spectateur est étudiant » et $A$ : « le spectateur a un abonnement ».
\begin{enumerate}[label=\alph*)]
  \item Décrire par une phrase les événements $E\cap A$, $E\cup A$, $\overline{E}$ et
    $E\cap\overline{A}$.
  \item Léo est étudiant et n'a pas d'abonnement. Lesquels des événements $E$, $A$,
    $E\cap A$ et $E\cup A$ sont réalisés ? Répondre par une phrase.
\end{enumerate}

\exo[Défauts de fabrication]
Dans la production d'une usine de téléphones, $4\,\%$ des appareils ont un défaut
d'écran, $6\,\%$ un défaut de batterie, et $1\,\%$ les deux défauts. On choisit un
téléphone au hasard ; on note $E$ : « il a un défaut d'écran » et $B$ : « il a un
défaut de batterie ».
\begin{enumerate}[label=\alph*)]
  \item Donner $P(E)$, $P(B)$ et $P(E\cap B)$.
  \item Calculer la probabilité que le téléphone n'ait pas de défaut d'écran.
  \item Calculer la probabilité qu'il ait au moins un des deux défauts.
  \item Calculer la probabilité qu'il n'ait aucun défaut.
\end{enumerate}

\exo[Des données impossibles]
$A$ et $B$ sont deux événements d'une même expérience. Pour chaque situation a), b),
c), expliquer par une phrase pourquoi elle est impossible.
\begin{enumerate}[label=\alph*)]
  \item $P(A)=0{,}7$ et $P\big(\overline{A}\big)=0{,}4$.
  \item $P(A)=0{,}6$ et $P(B)=0{,}5$, avec $A$ et $B$ incompatibles.
  \item $P(A)=0{,}3$ et $P(A\cap B)=0{,}45$.
  \item On sait maintenant que $P(A)=0{,}5$, $P(B)=0{,}35$ et $P(A\cup B)=0{,}7$.
    Calculer $P(A\cap B)$.
\end{enumerate}

\partiefiche{Loi de probabilité}

\exo[Un dé truqué]
On lance un dé truqué à six faces, dont la loi de probabilité est la suivante.
\begin{center}
\renewcommand{\arraystretch}{1.4}
\begin{tabular}{|l|c|c|c|c|c|c|}
\hline
Face & $1$ & $2$ & $3$ & $4$ & $5$ & $6$ \\
\hline
Probabilité & $0{,}15$ & $0{,}2$ & $0{,}1$ & $\dots$ & $0{,}25$ & $0{,}1$ \\
\hline
\end{tabular}
\end{center}
\begin{enumerate}[label=\alph*)]
  \item Calculer la probabilité d'obtenir $4$.
  \item Calculer les probabilités des événements $A$ : « obtenir un nombre pair » et
    $B$ : « obtenir au moins $4$ ».
  \item Décrire $\overline{B}$ par une phrase, puis calculer $P\big(\overline{B}\big)$.
\end{enumerate}

\exo[Une roue de quatre couleurs]
Une roue est partagée en secteurs rouges, verts, bleus et jaunes. La loi de la couleur
obtenue est la suivante.
\begin{center}
\renewcommand{\arraystretch}{2.4}
\begin{tabular}{|l|c|c|c|c|}
\hline
Couleur & rouge & vert & bleu & jaune \\
\hline
Probabilité & $\dfrac{1}{3}$ & $\dfrac{1}{4}$ & $\dfrac{1}{6}$ & $\dots$ \\
\hline
\end{tabular}
\end{center}
\begin{enumerate}[label=\alph*)]
  \item Calculer la probabilité d'obtenir jaune.
  \item Calculer la probabilité d'obtenir une couleur qui n'est ni rouge ni bleue.
  \item Les issues de cette expérience sont-elles équiprobables ? Justifier.
\end{enumerate}

\exo[Attente à la cantine]
Un grand nombre de relevés a donné la répartition du temps d'attente à la cantine
d'un lycée. On la prend comme loi de probabilité pour un jour quelconque.
\begin{center}
\renewcommand{\arraystretch}{1.4}
\begin{tabular}{|l|c|c|c|c|}
\hline
Attente (min) & $0$ à $5$ & $5$ à $10$ & $10$ à $15$ & plus de $15$ \\
\hline
Probabilité & $0{,}2$ & $0{,}4$ & $0{,}3$ & $\dots$ \\
\hline
\end{tabular}
\end{center}
\begin{enumerate}[label=\alph*)]
  \item Compléter le tableau.
  \item Calculer la probabilité d'attendre moins de $10$ minutes.
  \item Calculer la probabilité d'attendre plus de $5$ minutes, à l'aide de
    l'événement contraire.
\end{enumerate}

\exo[Une probabilité inconnue]
Un dé à six faces est truqué : les faces $1$ à $5$ ont la même probabilité $a$, et
la face $6$ a une probabilité trois fois plus grande.
\begin{enumerate}[label=\alph*)]
  \item Montrer que $8a=1$, puis calculer $a$.
  \item Dresser le tableau de la loi de probabilité.
  \item Calculer la probabilité d'obtenir un nombre pair.
  \item Calculer la probabilité de ne pas obtenir $6$.
\end{enumerate}

\partiefiche{Tableaux croisés d'effectifs}

\exo[Une salle de concert]
Une salle de concert a vendu $500$ places : $300$ en fosse, les autres en gradins.
$40\,\%$ des spectateurs en fosse et $15\,\%$ des spectateurs en gradins ont moins de
$25$ ans.
\begin{enumerate}[label=\alph*)]
  \item Dresser le tableau croisé des effectifs (lignes : fosse, gradins ;
    colonnes : moins de $25$ ans, $25$ ans ou plus ; totaux).
  \item On choisit un spectateur au hasard ; on note $F$ : « il est en fosse » et
    $J$ : « il a moins de $25$ ans ». Calculer $P(F)$, $P(J)$ et $P(F\cap J)$.
  \item Calculer $P(F\cup J)$, puis la probabilité que le spectateur soit en gradins
    et ait $25$ ans ou plus.
\end{enumerate}

\exo[Musique et dessin]
Dans une classe de $32$ élèves, $14$ font de la musique, $10$ font du dessin, et $4$
font les deux.
\begin{enumerate}[label=\alph*)]
  \item Dresser un tableau croisé des effectifs (musique ou non ; dessin ou non).
  \item Combien d'élèves ne font ni musique ni dessin ?
  \item On choisit un élève au hasard ; on note $M$ : « il fait de la musique » et
    $D$ : « il fait du dessin ». Décrire $M\cap\overline{D}$ par une phrase, puis
    calculer sa probabilité.
  \item Calculer $P(M\cup D)$.
\end{enumerate}

\exo[Satisfaction des usagers]
Un réseau de transports a interrogé $1\,000$ usagers : $400$ prennent le bus, $350$
le tram, les autres le métro. Sont satisfaits $70\,\%$ des usagers du bus, $80\,\%$
de ceux du tram et $90\,\%$ de ceux du métro.
\begin{enumerate}[label=\alph*)]
  \item Dresser le tableau croisé des effectifs (lignes : satisfait, non satisfait ;
    colonnes : bus, tram, métro ; totaux).
  \item On choisit un usager au hasard ; on note $T$ : « il prend le tram » et $S$ :
    « il est satisfait ». Calculer $P(T)$, $P(S)$ et $P(T\cap S)$.
  \item Calculer $P(T\cup S)$.
\end{enumerate}

\partiefiche{Probabilités conditionnelles}

\exo[Langue vivante et demi-pension]
Le tableau répartit les $500$ élèves de seconde d'un lycée. On note $D$ : « l'élève
est demi-pensionnaire », $E$ : « l'élève apprend l'espagnol » et $A$ : « l'élève
apprend l'allemand ».
\begin{center}
\renewcommand{\arraystretch}{1.4}
\begin{tabular}{|l|c|c|c|}
\hline
 & Espagnol & Allemand & Total \\
\hline
Demi-pensionnaires & $210$ & $90$ & $300$ \\
\hline
Externes & $170$ & $30$ & $200$ \\
\hline
Total & $380$ & $120$ & $500$ \\
\hline
\end{tabular}
\end{center}
\begin{enumerate}[label=\alph*)]
  \item On choisit au hasard un demi-pensionnaire. Calculer la probabilité qu'il
    apprenne l'allemand, et la noter.
  \item On choisit au hasard un élève qui apprend l'allemand. Calculer la probabilité
    qu'il soit demi-pensionnaire, et la noter.
  \item Calculer $P_E(D)$ (au centième près) et traduire cette probabilité par une
    phrase.
\end{enumerate}

\exo[De la phrase à la notation]
On choisit au hasard un élève d'un lycée ; on note $D$ : « l'élève est
demi-pensionnaire » et $S$ : « l'élève fait partie de l'AS ».
\begin{enumerate}[label=\alph*)]
  \item Traduire chaque phrase par une égalité : « $45\,\%$ des demi-pensionnaires
    font partie de l'AS » ; « parmi les élèves de l'AS, $60\,\%$ sont
    demi-pensionnaires » ; « $20\,\%$ des élèves sont des demi-pensionnaires de
    l'AS ».
  \item Traduire par une phrase : $P_{\overline{D}}(S)=0{,}3$.
\end{enumerate}

\exo[Un test de dépistage]
Un test de dépistage d'une maladie est essayé sur $20\,000$ personnes, dont $2\,\%$
sont malades. On note $M$ : « la personne est malade » et $T$ : « le test est
positif ».

Rappel : la \emph{sensibilité} est la part des malades dont le test est positif,
$P_M(T)$ ; la \emph{spécificité}, la part des non-malades dont le test est négatif,
$P_{\overline{M}}\big(\overline{T}\big)$ ; un \emph{faux positif} est un test positif
chez une personne non malade.

La sensibilité du test est $0{,}9$ et sa spécificité $0{,}97$.
\begin{enumerate}[label=\alph*)]
  \item Combien de personnes sont malades ? Combien d'entre elles ont un test
    positif ?
  \item Combien de personnes non malades ont un test négatif ? un test positif ?
  \item Dresser le tableau croisé des effectifs (malades ou non ; test positif ou
    négatif ; totaux).
  \item Une personne a un test positif. Calculer la probabilité qu'elle soit malade
    (au centième près).
  \item Quelle part des tests positifs sont des faux positifs ? Commenter.
\end{enumerate}

\exo[Avec la formule]
On tire au hasard une carte d'un jeu de $32$ cartes ($4$ couleurs de $8$ cartes : $7$,
$8$, $9$, $10$, valet, dame, roi, as). On note $C$ : « la carte est un cœur » et $R$ :
« la carte est un roi ».
\begin{enumerate}[label=\alph*)]
  \item Calculer $P(C)$, $P(R)$ et $P(C\cap R)$.
  \item Calculer $P_C(R)$ et $P_R(C)$ avec la propriété 6, puis traduire chacune par
    une phrase.
  \item Calculer $P(C\cup R)$.
\end{enumerate}

\partiefiche{Arbres pondérés}

\exo[Compétition au club de tennis]
Dans un club de tennis, $60\,\%$ des membres sont adultes. $30\,\%$ des adultes et
$70\,\%$ des jeunes jouent en compétition. On choisit un membre au hasard ; on note
$A$ : « il est adulte » et $C$ : « il joue en compétition ».
\begin{enumerate}[label=\alph*)]
  \item Traduire les données par des probabilités, puis construire l'arbre pondéré.
  \item Calculer $P(A\cap C)$.
  \item Calculer $P(C)$.
  \item Calculer la probabilité que le membre soit jeune et ne joue pas en
    compétition.
\end{enumerate}

\exo[Deux chocolats]
Une boîte contient $3$ chocolats noirs et $4$ chocolats au lait, indiscernables au
toucher. On en prend un au hasard, on le mange, puis on en prend un second.
\begin{enumerate}[label=\alph*)]
  \item Construire l'arbre pondéré.
  \item Calculer la probabilité de prendre deux chocolats noirs, puis celle de prendre
    deux chocolats au lait.
  \item Calculer la probabilité de prendre deux chocolats de la même sorte.
  \item Calculer la probabilité de prendre au moins un chocolat noir, à l'aide de
    l'événement contraire.
\end{enumerate}

\exo[Deux tickets à gratter]
Un ticket à gratter est gagnant avec la probabilité $\frac{1}{4}$. Max achète deux
tickets ; on admet que les résultats des deux tickets ne dépendent pas l'un de
l'autre.
\begin{enumerate}[label=\alph*)]
  \item Construire l'arbre pondéré (G : « le ticket est gagnant »).
  \item Calculer la probabilité que les deux tickets soient gagnants.
  \item Calculer la probabilité qu'un seul des deux tickets soit gagnant.
  \item Calculer la probabilité qu'au moins un ticket soit gagnant.
\end{enumerate}

\exo[Un test pour les chats]
Dans un refuge, $1\,\%$ des chats ont une maladie. Un test a une sensibilité de
$0{,}98$ (il est positif chez $98\,\%$ des chats malades) et une spécificité de
$0{,}95$ (il est négatif chez $95\,\%$ des chats non malades). On teste un chat
choisi au hasard ; on note $M$ : « le chat est malade » et $T$ : « le test est
positif ».
\begin{enumerate}[label=\alph*)]
  \item Construire l'arbre pondéré.
  \item Calculer $P(M\cap T)$, puis $P(T)$.
  \item Le test d'un chat est positif. Calculer la probabilité qu'il soit malade,
    $P_T(M)$, au centième près, avec la propriété 6.
  \item Comparer $P_M(T)$ et $P_T(M)$. Un test positif suffit-il pour conclure ?
\end{enumerate}

\partiefiche{Loi des grands nombres et simulation}

\exo[Lire une simulation]
On lance un dé équilibré à quatre faces numérotées de $1$ à $4$ ; on gagne si on
obtient $4$.
\begin{python}
from random import randint

def frequence_gain(n):
    gains = 0
    for k in range(n):
        de = randint(1, 4)
        if de == 4:
            gains = gains + 1
    return gains / n
\end{python}
\begin{center}
\renewcommand{\arraystretch}{1.4}
\begin{tabular}{|l|c|c|c|c|}
\hline
$n$ & $10$ & $100$ & $1\,000$ & $100\,000$ \\
\hline
Résultat & $0{,}4$ & $0{,}21$ & $0{,}262$ & $0{,}24937$ \\
\hline
\end{tabular}
\end{center}
\begin{enumerate}[label=\alph*)]
  \item Calculer la probabilité $p$ de gagner.
  \item Expliquer le rôle de la variable \texttt{gains}, puis ce que renvoie la
    fonction.
  \item Que montrent les résultats du tableau ?
  \item Pour $n=100$, puis pour $n=1\,000$, l'écart entre la fréquence et $p$ est-il
    inférieur à $\dfrac{1}{\sqrt{n}}$ ?
\end{enumerate}

\exo[Simuler un contrôle de qualité]
Dans une imprimerie, $4\,\%$ des livres ont un défaut. On simule le contrôle de livres
pris au hasard dans un très grand stock.
\begin{python}
from random import random

def livre():
    if random() < .....:
        return 1
    else:
        return .....

def frequence(n):
    defauts = 0
    for k in range(n):
        defauts = defauts + livre()
    return .....
\end{python}
\begin{enumerate}[label=\alph*)]
  \item Compléter les fonctions : \texttt{livre()} renvoie $1$ si le livre a un
    défaut, $0$ sinon ; \texttt{frequence(n)} renvoie la fréquence des livres
    défectueux sur $n$ livres.
  \item De quelle valeur \texttt{frequence(10000)} sera-t-il proche ?
  \item Dans environ $95\,\%$ des cas, entre quelles valeurs se trouve cette
    fréquence pour $n=10\,000$ ?
\end{enumerate}

\exo[Estimer une probabilité]
On lance $2\,500$ fois une punaise : elle tombe $1\,550$ fois sur la tête.
\begin{enumerate}[label=\alph*)]
  \item Calculer la fréquence observée $f$ de l'issue « sur la tête ».
  \item Calculer $\dfrac{1}{\sqrt{2\,500}}$. Entre quelles valeurs se trouve
    vraisemblablement la probabilité $p$ que la punaise tombe sur la tête ?
  \item Combien de lancers faudrait-il pour que l'écart entre $f$ et $p$ soit
    inférieur à $0{,}01$ dans environ $95\,\%$ des cas ?
\end{enumerate}

\end{multicols}

\end{document}
